\documentclass{leveldoc}
% Level @NUMBER@ (@TITLE@): the physics of this level, derived from first principles.
% Written to docs/levels/README.md and docs/levels/GROK.md; the examples are
% docs/levels/first_bend (a closed form) and docs/levels/slingshot (a numerical solution).
% Every number comes from calc.py (\val{...}); every comment starting with "TEMPLATE" is
% replaced or deleted before the document is done.
%
% \DocUses: the documents cited by equation or section; always foundations, then earlier
% levels only (e.g. \DocUses{foundations,first_bend}).
\DocUses{foundations}
% The level file's fingerprint the text was last checked against: set by
% `python scripts/level_docs.py new`; when the level changes, `check` warns and prints the
% new value, to be set after the text has been reviewed.
\DocFingerprint{@FINGERPRINT@}

\begin{document}

\section{What the level asks}\label{sec:asks}

The game's text:
\leveldescription

% TEMPLATE. The puzzle in physical terms: what is launched, from where, where it must go,
% what the player controls (the limits), and what ends a flight
% (\xsecref{foundations}{sec:board}). State the goal exactly (boxes, acceptance windows,
% goal values and their time windows). Name the reference solution and say in one sentence
% why it works. Say what this document does: a closed form, or a numerical solution.

\section{Setup}\label{sec:setup}

The level file, field by field (the appendix reproduces it):
\setuptable

% TEMPLATE. The symbols for every quantity the derivation uses, each with its value from
% calc.py (\val{...}), and the geometry (positions, boxes, the world, t_max). A figure of the
% board (TikZ/pgfplots; coordinates in cells; data from calc.py with doc.data and
% \addplot table {name.dat}) helps; use \rawval{...} for numbers inside TikZ coordinates.

\section{Principles and model}\label{sec:model}

\begin{principles}
  % TEMPLATE. Each principle with its source: an equation of the Foundations
  % (\xeqref{foundations}{eq:...}), of an earlier level (\xeqref{slug}{eq:...}), or a
  % textbook (\Jackson{§...}, \LandauII{§...}, ...: cite a section only where checked).
  \item
\end{principles}

% TEMPLATE. The model as equations: the equations of motion (or of the fields) specialised to
% this level, with their initial and boundary conditions and the rule that ends the flight.
% Then every assumption or approximation, each in its own box with what it neglects and how
% large that is (a number from calc.py). Approximations the game makes too, and the ones only
% this document makes, are both listed; say which is which.

\begin{assumption}{TEMPLATE: a short name}\label{as:first}
TEMPLATE. What is assumed, why it holds, and the size of what it neglects.
\end{assumption}

\section{Derivation}\label{sec:derivation}

% TEMPLATE. From the principles to formulas that can be evaluated: every step shown or cited
% (no "it can be shown"). Subsections as needed. Number every equation that is used later
% (\label{eq:...}). Where an earlier document derived something, cite it instead of
% repeating it. Where no closed form exists, state the exact mathematical problem and the
% numerical method that solves it (the Foundations' toolkit: \xsecref{foundations}{sec:numerics})
% with its error estimate. Where a conserved quantity or a limit gives an independent check,
% derive it.

\section{The reference solution by hand}\label{sec:hand}

% TEMPLATE. The formulas evaluated for the reference solution, as numbered steps a student
% can follow with a calculator (and, where needed, a stated numerical method run for long
% enough). Every intermediate value from calc.py. End with the result.

\begin{procedure}
  \item\label{st:first} TEMPLATE.
\end{procedure}

\begin{result}[The reference flight]
TEMPLATE. Where and when the level's goal is met, with the numbers.
\end{result}

\section{Comparison with the game}\label{sec:comparison}

\comparisontable

% TEMPLATE. Discuss the differences: are they the size the error estimates predict (the
% game's verified accuracy, \xsecref{foundations}{sec:verification}; this document's own
% method and approximations)? The tolerances in calc.py were set from those estimates before
% comparing; never widen one to make a row pass.

\section{Other solutions}\label{sec:others}

% TEMPLATE. Why the reference works and what else does: the other placements that solve the
% level, a map of them if it is cheap to compute (a family in placements.json, see
% docs/levels/first_bend), and a few alternatives checked in detail against the game
% (placements.json "detailed"). Where the level's sliders have a window, explain its edges.

\section{What the game computes}\label{sec:engine}

% TEMPLATE. The engine's method for this level (integrator, field solver, events, verification)
% with references to PHYSICS.md (\physref{...}), what in it is exact and what approximate, and
% why its numbers and this document's agree (or how far they may differ).

\section{Further reading}\label{sec:reading}

\begin{itemize}
  % TEMPLATE. Textbook sections on this level's physics (checked), and earlier levels that
  % build up to it (\levelref{slug}).
  \item
\end{itemize}

\levelfileappendix

\end{document}
